\section{Partial effects and recovery planning}
\label{sec:C10:recovery}
External operations can partially succeed. A file edit may commit before a test process crashes. A remote request may complete before a timeout. A multi-step workflow may create one resource and fail while linking another. Recovery begins by identifying the actual state, not by assuming the last returned error describes every effect. The planner should use operation identifiers, status queries, and recorded postconditions where supported. Blindly repeating a non-idempotent operation can turn an uncertain outcome into a definite duplicate.

Compensation is not the same as reversal. Deleting a newly created record may undo some state while leaving an email, audit event, charge, or downstream observation irreversible. A recovery plan should state which effects can be restored and which require disclosure or further action. This distinction constrains autonomous authority: a system may safely retry a read but need stronger authorization for an operation whose compensation is incomplete. A model's confidence that it can fix a mistake later is not evidence that the environment supports reversal.

The plan should reserve resources for verification and recovery. Spending the entire budget on candidate generation can leave no capacity to inspect the final outcome. A deadline-aware controller can allocate a minimum final-check budget and stop exploration earlier. This is a constrained scheduling problem coupled to inference-time control in Chapter~\ref{ch:C09}. Tool latency distributions matter: an operation with a short mean but a heavy tail can consume the recovery window. Report timeout and censored outcomes rather than treating them as if the agent had simply chosen to stop.

Fault injection can test recovery without inventing real failures in production. In a controlled environment, interrupt an operation after a known partial effect, drop a response, revoke permission between planning and execution, or replace a stored checkpoint with a stale one. Verify whether the agent reconstructs state, avoids duplicate effects, and preserves unresolved uncertainty. Keep these tests isolated and reversible. They diagnose mechanisms and do not estimate deployment incident frequencies unless the injected distribution is itself justified by operational data.

